\section{Introduction}
\label{sec:introduction}

An autonomous agent is asked to post a journal entry, issue a refund, or file a
regulatory notification. It calls an enterprise API written years before any of
this. The call times out, and nothing on either side can say whether the money
moved. Autonomous agents are the setting in which this is now most visible, and
we use them throughout as the motivating example. The problem itself belongs to
the endpoint, and any caller that can fail mid-call inherits it. A workflow
engine, a retrying batch job and a human operator with a script are all in the
same position.

This is an old problem in a new setting, and the new setting removes the
assumption that the old solutions rest on. Durable-execution engines, workflow
systems and message brokers achieve their guarantees by requiring that the
operation being retried is idempotent, or that the remote side will accept a
deduplication key, or that both endpoints can be enrolled in one
transaction~\cite{helland2007apostate,burckhardt2021durable,zhuang2023exoflow}.
Temporal's documentation is explicit about the precondition: ``You should
always make your business logic Activities idempotent in Temporal. Because
Activities may be retried, these functions may be executed more than
once''~\cite{temporal-activity-definition}. The legacy endpoints at issue satisfy
none of these conditions. They are non-idempotent and they accept no idempotency
key. Many of them also cannot be asked afterwards whether a mutation was
applied, and that third gap is the one that does the real damage.

\subsection{The trilemma}
\label{sec:trilemma}

Consider a worker that has sent a mutation and then crashed, and a recovery
process that must decide what to do. Exactly three things can happen to the
external world, and a system must choose which one it is willing to produce.

\begin{itemize}
\item An \textbf{\undetdup{}}. The recovery process re-sends. If the first
      request had in fact been applied, a second real effect now exists and no
      component of the system knows it. A naive retry produces this outcome,
      and so, as \cref{sec:evaluation} shows, do a lease, a fenced
      compare-and-swap and a durable event-sourced history, which each fail
      to prevent it because none of them is a record written before the
      call.
\item A \textbf{\lostfx{}}. The recovery process does not re-send and records
      the attempt as failed. If the first request was applied, a real effect
      now exists that the system's records deny. An at-most-once configuration
      produces this outcome. It is quieter than a duplicate and no better: a
      refund that the ledger says did not happen becomes a reconciliation
      problem that surfaces in the accounts weeks later.
\item A \textbf{\declamb{}}. The system declines to decide, records that it
      cannot decide, and stops. The effect may or may not exist. What the
      system guarantees is that a durable record says so, and that no further
      automated action is taken on that execution.
\end{itemize}

The first two outcomes are silent and the third is not. Our position is that
the third corner is the one a protocol should be engineered to reach. Against
an endpoint that cannot be queried, avoiding ambiguity altogether is
impossible, so the quantity a system should be judged on is whether its
residual uncertainty is declared and bounded rather than silent and unbounded.

This impossibility is not incidental. Deciding whether a remote effect occurred,
when the channel may fail and the remote side may not answer, is the
coordinated-attack problem; no protocol resolves it with certainty, and the
classical results on consensus under crash faults and on failure detection say
why~\cite{fischer1985impossibility,chandra1996failuredetectors}. What a designer
can choose is where the uncertainty is allowed to surface. We call our protocol
AEP. It records an intent before every external call and fails closed when
recovery cannot resolve the outcome. AEP places the uncertainty in a durable
state that an operator can see, and in every cell we measured, none of it
reached the accounts.

\subsection{Contributions}

\begin{itemize}
\item[\textbf{C1}] \textbf{The declared-ambiguity formulation.} We state the
      problem as the three-way trade of \cref{sec:trilemma} rather than as a
      pursuit of exactly-once. The formulation gives three properties: fenced
      state, detectable ambiguity, and a fail-closed liveness bound. Each is
      mapped to the code path that enforces it, and each has its residual
      window declared rather than assumed away (\cref{sec:model,sec:protocol}).
\item[\textbf{C2}] \textbf{A protocol and an implementation} in which a
      durably acknowledged write-ahead intent is a checked precondition of
      dispatch authority, not a matter of call ordering. Under a trusted-code
      assumption, the fsync acknowledgment returns an opaque, non-copyable,
      single-use, scope-bound in-process dispatch guard. The guard is consumed
      to write a Redis-visible authorization, and the pre-dispatch script
      refuses to proceed without it (\cref{sec:protocol}). This is a
      control-flow invariant of the supported API. It is not a cryptographic
      capability, and it is not a boundary against arbitrary code in the same
      Python process.
\item[\textbf{C3}] \textbf{An evaluation under real process kills} across six
      named crash points, three endpoint reconciliation capabilities and
      five reported fault regimes, against five baseline designs, one of them
      in two configurations. The five regimes are every execution crashed,
      none crashed, Redis hard-killed before the acknowledgment, Redis killed
      after dispatch, and storage discarding writes while reporting success.
      (A sixth regime, a $30\%$ crash probability, was collected; its rates
      are not in these tables and the collection is in the artifact.)
      % Regimes counted from each collection root's analysis/coverage.json.
      % The five reported: crashed and p0 from experiments/results/matrix,
      % redis-kill-preack from matrix and the b2/phase13-armA sessions,
      % redis-kill-inflight from phase13-inflight-s1,s2 (6.3.2.5), and
      % write-loss-preack from reports/raw/ws4-writeloss-s1-2026-09-07
      % (6.3.4). The sixth is p30, ws5-2026-09-10/t2-p30, which no table
      % quotes. (session-3) and crashed are one regime under the two labels
      % docs/29 3 declares as a normalisation, so they are counted once.
      % This said 'three collected fault regimes' until 2026-09-17, which
      % was the matrix root's three and excluded the two regimes 6.3.2.5
      % and 6.3.4 report -- one of them the fault the abstract's prevention
      % sentence turns on.
      AEP records no
      undetected duplicate and no lost effect in any cell measured, and the
      baselines without a pre-dispatch record duplicate in most crashed
      executions. AEP's residual is declared ambiguity, whose rate is a
      function of what the endpoint can be asked (\cref{sec:evaluation}).
\item[\textbf{C4}] \textbf{A decomposition of the mechanism, by ablation,
      that reassigns our own headline result.} The write-ahead pattern is
      normally presented as one mechanism delivering one guarantee, and it is
      two. The pre-dispatch record together with no re-entry
      is what makes an outcome detectable, and it does so without the
      durability barrier: ablating the barrier produces no observed difference
      in the crashed-regime detection metrics (\cref{sec:eval-detection}). The
      barrier contributes something else, withholding the external call when
      durability can no longer be confirmed. That contribution is invisible to
      the crash faults that dominate this literature and large under the fault
      that targets Redis (\cref{sec:eval-prevention}). Separating the two
      changes what a deployment must pay for: detection is nearly free,
      prevention is where the fsync cost lives, and an operator can buy the
      first without the second (the supplementary material, \emph{The barrier
      is a deployment choice}).
\end{itemize}

\subsection{What this paper does not claim}

We claim no exactly-once external execution, no prevention of a duplicate the
provider has already accepted, no high availability, consensus or split-brain
immunity, and no durability beyond one local AOF fsync. The system is a single
Redis trust domain and the measurements are single-node. These are stated as a
table of non-claims in \cref{sec:model} with the reason each is out of scope,
and revisited in \cref{sec:threats}.
